% ABOUTME: Base LaTeX template for research-report plugin reports.
% ABOUTME: Argument-driven structure: front Synthesis (TL;DR), Background, dynamic body chapters with argument-style titles, back Conclusions, Open Questions, Methodology appendix.

\documentclass[11pt]{article}
\usepackage[margin=1.15in]{geometry}
\usepackage{hyperref}
\usepackage{booktabs}
\usepackage{enumitem}
\usepackage{natbib}
\usepackage{graphicx}
\usepackage{longtable}
\usepackage{parskip}
\usepackage{setspace}
\usepackage{titlesec}

% Paragraph and spacing settings for readability
\setlength{\parskip}{1.0em}
\setlength{\parindent}{0pt}
\setstretch{1.35}

% Tighter section spacing
\titlespacing*{\section}{0pt}{1.5em}{0.8em}
\titlespacing*{\subsection}{0pt}{1.2em}{0.6em}
\titlespacing*{\subsubsection}{0pt}{1em}{0.4em}

% Compact list settings
\setlist[itemize]{itemsep=0.3em, parsep=0.2em, topsep=0.5em}
\setlist[enumerate]{itemsep=0.3em, parsep=0.2em, topsep=0.5em}

\hypersetup{
  colorlinks=true,
  linkcolor=blue,
  citecolor=blue,
  urlcolor=blue
}

\title{PLACEHOLDER_TITLE}
\author{Research Report --- Generated via Autonomous Workflow}
\date{\today}

\begin{document}
\maketitle
\tableofcontents
\newpage

%% ==========================================================================
%% WRITING PRINCIPLES (read these before writing ANY content):
%%
%% STRUCTURAL
%% 1. The body is built from chapter-level ARGUMENTS, not topic buckets.
%%    Every chapter heading should be a sentence that takes a position.
%%    BAD: "The State of LLM Agents"
%%    GOOD: "LLM agents have not crossed the reliability threshold for
%%           production autonomy"
%% 2. Each body chapter is a self-contained piece of argumentative writing
%%    that defends its chapter heading using cited evidence from the pool.
%% 3. Single authorial voice across the whole report. The narrative-writer
%%    agent reads voice-guide.md before each chapter and follows it.
%%
%% PROSE
%% 4. Every paragraph should make a CLAIM and use \cite{} as evidence FOR
%%    that claim. Forbid paragraphs that are just sequential cite-statements
%%    ("Source X says Y. Source Z says W.") — that is notes, not writing.
%% 5. Each section/chapter must close with a paragraph that INTERPRETS,
%%    not describes. State the judgment, the so-what, the upshot. Form is
%%    flexible — an unlabeled closing paragraph or a labeled
%%    \paragraph{Implications} when the so-what is especially load-bearing.
%% 6. Maximum 4-5 sentences per paragraph. Break longer ones up.
%% 7. Use \begin{itemize}/\begin{enumerate} for ANY list of 3+ items.
%%    Do NOT inline lists as comma-separated items in prose.
%% 8. Use \subsubsection{} to break up long chapters into scannable chunks.
%%
%% RIGOR
%% 9. Every factual claim has an inline \cite{key}.
%% 10. Preserve qualifications and evidence-gap ratings from the pool.
%%     If a pool entry says "supports the prevalence claim but not the
%%     causal claim", the prose must reflect that boundary.
%% 11. Reader-pass agents may improve flow but MUST NOT drop \cite{},
%%     weaken qualifications, or turn qualified claims into unqualified ones.
%% ==========================================================================

\section{Synthesis}
% NOTE: This is the front-loaded executive summary. Written in Phase S
% AFTER the body chapters and Conclusions are finalized, so it accurately
% reflects what was actually argued.
%
% Purpose: a reader who reads only this section understands the core
% findings, the argument the report makes, and the action items.
% Length: 1500-2000 words. Absolute max 2500.
% Self-contained, no temporal narrative, no references to iterations.

\subsection{Summary}
% 2-3 short paragraphs (3-4 sentences each). What was researched, why it
% matters, scope, the report's headline argument in one sentence.

\subsection{Key Takeaways}
% 5-7 most important findings as a numbered list (\begin{enumerate}).
% Each item: \item \textbf{Takeaway title.} 2-4 sentences of explanation
% with \cite{} references. Each must reference its supporting chapter
% (e.g., "see Chapter 3").
% Takeaways must be internally consistent — no contradictions between items.

\subsection{Conclusions \& Recommendations Preview}
% Brief 2-3 sentence pointer to the back-loaded Conclusions section.
% The full conclusions are at the end of the report; this is just a
% signpost so a TL;DR reader knows where to find the action items.

\subsection{Confidence \& Limitations}
% One short paragraph for overall confidence, then \begin{itemize} listing
% specific limitations (1-2 sentences each).

\section{Background \& Context}
% Why this question matters, what motivated the research, the prior state
% of understanding, key terms used in the report.
% Keep paragraphs short (3-5 sentences). Use itemize for listing context items.
% Define key terminology once — voice-guide.md records the terminology
% decisions and the rest of the report follows them consistently.

%% ==========================================================================
%% BODY CHAPTERS
%%
%% This region is populated dynamically in Phase S by the narrative-writer
%% agent. Each chapter:
%%   - Has an argument-style \section{} heading (a sentence that takes a
%%     position, not a topic bucket).
%%   - Defends that argument using cited evidence from the pool.
%%   - Closes with an interpretive paragraph stating the so-what.
%%
%% Chapter count: typically 3-7, capped at 8. Long topics use deeper
%% chapters, not more chapters.
%%
%% The narrative-writer drafts one chapter at a time, sequentially, so a
%% single voice carries through. Do NOT parallelize chapter drafting.
%%
%% Example (placeholder — replace during Phase S):
%%
%% \section{[Chapter 1's argument as a sentence taking a position]}
%% [Continuous argumentative prose with inline \cite{} references.]
%% [Use \subsection{} for major sub-arguments within the chapter.]
%% [Closing paragraph or \paragraph{Implications} with the so-what.]
%%
%% \section{[Chapter 2's argument...]}
%% ...
%% ==========================================================================

% CHAPTERS_PLACEHOLDER — narrative-writer inserts dynamic body here.

\section{Conclusions \& Recommendations}
% NOTE: This is the back-loaded Conclusions. Written in Phase S AFTER all
% body chapters are drafted. The integrated argument earned by everything
% in between.
%
% Purpose: state the unified position the body chapters jointly support,
% and the actions that follow from it. The reader has now done the work,
% so this section can be more direct than the front Synthesis.
% Length: 1000-2000 words.

\subsection{The Argument, Stated Plainly}
% One paragraph stating the report's overall thesis as it now stands,
% having walked through the evidence. Reference the chapters that earned it.

\subsection{Conclusions}
% Use \begin{itemize} or numbered prose. Each conclusion is a position
% the report defends, with the chapter(s) that support it cited explicitly.
% Be opinionated. Conclusions are judgments, not summaries.

\subsection{Recommendations}
% Use \begin{itemize}. Each recommendation is paired with the conclusion
% it follows from. Specify: who should do it, why now, what changes if
% they don't. Recommendations must flow logically from their paired
% conclusions — no orphan recommendations.

\subsection{Conditions Under Which the Argument Could Change}
% Brief — what evidence or development would warrant revising the
% conclusions? This is the intellectual-honesty appendix to the
% conclusions, distinct from Open Questions.

\section{Open Questions}
% Use enumerate or itemize — this is naturally a list.
% Each question: 1-2 sentences explaining why it remains open and what
% would be needed to settle it.

\appendix

\section{Methodology}
% Iterations run, sources consulted, source credibility hierarchy used,
% how chapter-level arguments evolved during research, key revisions to
% the thesis, evidence-pool size, reader-pass count.
% This is the audit trail. Use itemize and tables liberally.

\bibliographystyle{plainnat}
\bibliography{sources}

\end{document}
